\documentclass[11pt]{article}

\usepackage{amsmath,amssymb,amsthm}
\usepackage{booktabs}
\usepackage{tabularx}
\usepackage{geometry}
\usepackage{hyperref}
\usepackage{xcolor}

\makeatletter
\g@addto@macro\UrlBreaks{\do\_\do\.}
\makeatother

\geometry{margin=2.7cm}
\setlength{\emergencystretch}{3em}
\hypersetup{
  colorlinks=true,
  linkcolor=blue!60!black,
  citecolor=green!45!black,
  urlcolor=blue!60!black,
  pdftitle={Erdos Problem 848: A Kernel-Checked Proof of the Exact Extremal Bound},
  pdfauthor={Alex Chengyu Li}
}

\newtheorem{theorem}{Theorem}[section]
\newtheorem{proposition}[theorem]{Proposition}
\newtheorem{lemma}[theorem]{Lemma}
\newtheorem{corollary}[theorem]{Corollary}
\theoremstyle{definition}
\newtheorem{definition}[theorem]{Definition}

\newcommand{\NN}{\mathbb{N}}
\newcommand{\CN}{C_N}
\newcommand{\DN}{D_N}
\newcommand{\TN}{T_N}
\newcommand{\RN}{R_N}
\newcommand{\leanname}[1]{}
\newcommand{\FiveMillion}{5\cdot10^6}
\newcommand{\TenMillion}{10^7}
\newcommand{\TwentyMillion}{2\cdot10^7}
\newcommand{\FortyMillion}{4\cdot10^7}
\newcommand{\TwoHundredMillion}{2\cdot10^8}
\newcommand{\TwoBillion}{2\cdot10^9}
\newcommand{\FiveHundredBillion}{5\cdot10^{11}}
\newcolumntype{Y}{>{\raggedright\arraybackslash}X}

\title{\textbf{Erd\H{o}s Problem 848:\\
A Kernel-Checked Proof of the Exact Extremal Bound}}
\author{Alex Chengyu Li\\[3pt]
\small Independent Researcher, Tokyo, Japan}
\date{August 2026}

\begin{document}
\maketitle

\begin{abstract}
For $N\geq1$, let $A\subseteq[1,N]$ satisfy the condition that $ab+1$ is
nonsquarefree for every $a,b\in A$, with $a=b$ allowed.  We determine the
exact maximum cardinality of such a set, proving that
\[
 |A|\leq \left\lfloor\frac{N+18}{25}\right\rfloor,
\]
with equality attained by the elements of the progression $7\pmod {25}$ in
$[1,N]$.  This establishes the extremal value for every $N$ and closes the
finite range left by the preceding asymptotic results.  The proof begins
with an exact Hall reformulation.  Exact prefix colourings handle
$N\leq 5\cdot10^6$.  For larger $N$, square-divisor estimates match the
second mod-$25$ construction and force any remaining Hall defect to contain
many elements outside both distinguished progressions.  We partition those
elements by their five possible $2$-adic types and their residues modulo $9$.
Three- and four-pivot arguments then bound the Hall completion through
explicit finite-prime counts and the solutions of a transformed quadratic
equation.  The remaining comparisons reduce to exact rational inequalities
over exhaustive finite cases, followed by a uniform estimate for the
unbounded tail.  Finite witnesses enter Lean as ordinary data, semantic
checkers establish the required inequalities, and the Lean~4 kernel replays
the final all-$N$ theorem with \texttt{--trust=0}.
\end{abstract}

% Machine-readable endpoint tags.  These macros are typographically empty.
\leanname{NonSquarefreeProductProp}
\leanname{OriginalProblem848Statement}
\leanname{originalA7_has_property}
\leanname{originalProblem_of_hallStatement}
\leanname{erdos848HallStatement_iff_originalProblem}
\leanname{erdos848_through_five_million}
\leanname{PaperCertificateProvider.fiveToTenMillion}
\leanname{PaperCertificateProvider.tenToTwentyMillion}
\leanname{PaperCertificateProvider.twentyToFortyMillion}
\leanname{erdos848FortyMillionClose_kernel}
\leanname{PaperCertificateProvider.twoHundredToTwoBillion}
\leanname{PaperCertificateProvider.twoBillionTail}
\leanname{PaperCertificateProvider.fortyMillionTail}
\leanname{erdos848_paper_tail_close}
\leanname{PaperGeneratedCertificateProvider.all_N}

\section{Introduction}

Write
\[
  \CN=\{n\in[1,N]:n\equiv7\pmod {25}\}.
\]
The quantifiers in the problem include the diagonal case $a=b$.  Moreover,
\begin{equation}\label{eq:sharp-cardinality}
 |\CN|=\left\lfloor\frac{N+18}{25}\right\rfloor.
\end{equation}
Erd\H{o}s and S\'ark\"ozy asked whether this construction is extremal
\cite{Erdos848Original}.  Sawhney proved the assertion for all sufficiently
large $N$ by a density sieve and structural reduction
\cite{Sawhney848}.  Sothanaphan made the asymptotic inputs explicit and
obtained the threshold $N\geq2.64\cdot10^{17}$
\cite{Sothanaphan848}.  We prove the exact statement
\begin{equation}\label{eq:target}
  |A|\leq|\CN|
\end{equation}
for every positive integer $N$.

\begin{table}[ht]
\centering
\footnotesize
\renewcommand{\arraystretch}{1.16}
\begin{tabularx}{\textwidth}{@{}>{\raggedright\arraybackslash}p{0.15\textwidth}Y Y Y@{}}
\toprule
Work & Proven range and conclusion & Remaining obstruction & Argument used here\\
\midrule
Sawhney (2025) &
The class $7\bmod25$ is extremal for every $N\geq N_0$, for some
unspecified $N_0$; a stability statement identifies the two candidate
extremal classes. &
The argument is asymptotic and gives neither $N_0$ nor the finite prefix. &
Density sieves and the mod-$25$ structural reduction.\\
\addlinespace
Sothanaphan (2026) &
The same extremal bound holds explicitly for
$N\geq2.64\cdot10^{17}$. &
The interval below that threshold remains; a pointwise search would not be a
practical proof. &
Explicit squarefree-progression estimates and a Pell-equation bound for
large square divisors.\\
\addlinespace
This paper &
The exact extremal value is proved for every $N\geq1$. &
No range remains. &
 Hall reformulation, interval estimates, valuation-class case analysis,
quadratic root counts, and exact rational inequalities.\\
\bottomrule
\end{tabularx}
\caption{Prior results and the all-$N$ theorem.}
\label{tab:prior-work}
\end{table}

The proof does not enumerate the possible values of $N$.  The Hall
reformulation replaces the extremal-set question by one inequality for a
completion $H_N(B)$.  Monotonicity permits a long interval of $N$ to be
controlled at its lower endpoint once its structural case is fixed.  The
cases arise from the $2$-adic class of a pivot, its residue modulo $9$, and
the quadratic-residue restrictions imposed on the complementary quotient in
$bx+1=p^2m$.  Exact finite calculations enter only after these reductions;
they prove rational inequalities and do not extrapolate from sampled values
\cite{HelfgottPlatt,OliveiraSilvaHerzogPardi}.

The squarefree-progression estimates belong to the classical line of
Prachar and Hooley \cite{PracharSquarefreeAP,HooleySquarefreeAP}.  Here they
are combined with exact residue and valuation information tailored to the
Hall graph.  Finite tables are accompanied by general soundness theorems, as
in other formalized proofs with substantial computation
\cite{GonthierOddOrder,HalesFlyspeck}.

\begin{proposition}[Sharpness]\label{prop:sharpness}
The set $\CN$ is admissible: if $a,b\in\CN$, then $ab+1$ is not squarefree.
\end{proposition}

\begin{proof}
Since $a\equiv b\equiv7\pmod {25}$,
\[
  ab+1\equiv 7^2+1=50\equiv0\pmod {25}.
\]
Thus $5^2\mid ab+1$.
\end{proof}

\section{The exact Hall cut}

Throughout, $N$ is a positive integer and $[1,N]=\{1,\ldots,N\}$.  An integer
is \emph{nonsquarefree} if it is divisible by $p^2$ for some prime $p$.
Join $b\in[1,N]\setminus\CN$ to $a\in\CN$ when $ab+1$ is squarefree.
For $B\subseteq[1,N]\setminus\CN$, put
\[
  \Gamma_N(B)=\{a\in\CN:\text{$ab+1$ is squarefree for some $b\in B$}\}
\]
and
\[
  \TN(B)=\CN\setminus\Gamma_N(B).
\]
Thus $a\in\TN(B)$ exactly when $ab+1$ is nonsquarefree for every $b\in B$.
Call $B$ a \emph{compatible outside set} if
\[
 B\subseteq[1,N]\setminus\CN
 \quad\text{and}\quad
 bb'+1\text{ is nonsquarefree for every }b,b'\in B,
\]
where $b=b'$ is allowed.  This is the exact condition denoted by
\texttt{Erdos848OutsideSet} together with
\texttt{NonSquarefreeProductProp} in the Lean development.

The usual Hall condition for the squarefree-edge graph is
$|\Gamma_N(B)|\geq |B|$.  Since
$|\Gamma_N(B)|=|\CN|-|\TN(B)|$, it is equivalent to
\begin{equation}\label{eq:hall-defect}
  |B|+|\TN(B)|\leq|\CN|.
\end{equation}
Only compatible outside sets need to be tested: they are precisely the
possible outside portions of an admissible set.

\begin{theorem}[Hall reduction]\label{thm:hall-reduction}
If \eqref{eq:hall-defect} holds for every compatible outside set $B$, then
\eqref{eq:target} holds.
\end{theorem}

\begin{proof}
Let $A$ be admissible and put $B=A\setminus\CN$.  Then
$A\cap\CN\subseteq\TN(B)$, whence
\[
 |A|=|B|+|A\cap\CN|
 \leq |B|+|\TN(B)|\leq|\CN|.
\]
\end{proof}

\begin{theorem}[Hall equivalence]\label{thm:hall-equivalence}
The Hall inequality \eqref{eq:hall-defect} is equivalent to the exact
upper bound \eqref{eq:target}.
\end{theorem}

\begin{proof}
The forward implication is Theorem~\ref{thm:hall-reduction}.  Conversely,
assume \eqref{eq:target} and let $B$ satisfy the hypotheses of
\eqref{eq:hall-defect}.  Put $A=B\cup\TN(B)$.  Products internal to $B$ are
nonsquarefree by hypothesis; products internal to $\TN(B)$ are divisible by
$25$; and every cross-product is nonsquarefree by the definition of
$\TN(B)$.  Hence $A$ is admissible.  The two parts are disjoint, so
\[
 |B|+|\TN(B)|=|A|\leq|\CN|.
\]
\end{proof}

For later normalization define the Hall completion
\[
  H_N(B)=B\cup\TN(B).
\]
The union is disjoint, so $|H_N(B)|=|B|+|\TN(B)|$.  The exact lower bound
\begin{equation}\label{eq:normalized-target}
  |\CN|\geq \frac N{25}-\frac7{25}
\end{equation}
shows that
\[
 \frac{|H_N(B)|}{N}\leq \frac1{25}-\frac7{25N}
\]
implies \eqref{eq:hall-defect}.  All rational comparisons below are made
against this target.
We call a compatible outside set with
$|\CN|<|B|+|\TN(B)|$ a \emph{strict Hall defect}.

\section{The range up to five million}

\begin{theorem}[Exact prefix colouring]
\label{thm:finite-prefix-colouring-close}
For every $1\leq N\leq\FiveMillion$, let
\[
 \mathcal D_N=\{x\in[1,N]:x^2+1\text{ is nonsquarefree}\},
\]
and join distinct $x,y\in\mathcal D_N$ when $xy+1$ is nonsquarefree.
This graph has a proper colouring whose colour set is indexed by $\CN$.
Consequently \eqref{eq:target} and \eqref{eq:hall-defect} hold throughout
this interval.
\end{theorem}

\begin{proof}
The diagonal hypothesis shows that every member of an admissible set belongs
to $\mathcal D_N$, and the off-diagonal hypothesis makes the admissible set a
 clique in this graph.  Thus a proper colouring by at most $|\CN|$ colours
 implies the cardinality bound.  Here ``indexed by $\CN$'' means a map into
 $\CN$; some colours may be unused.  A single finite certificate encodes the
 colourings for all prefixes in this range.

Put $L=\FiveMillion$.  At a prefix write
\[
 U=\{u\leq N:u\equiv7\pmod {25}\},\qquad
 V=\{v\leq N:v\equiv18\pmod {25}\},\qquad
 E=\mathcal D_N\setminus(U\cup V).
\]
For each $u\equiv7\pmod {25}$ and each $N\in[u,L]$, a state has three
possible occupants: the anchor $u$, at most one element of $V$, and at most
one element of $E$.  Empty non-anchor positions are recorded by zero.  The
 state is required to satisfy that every two positive occupants in different positions have
squarefree product plus one.  Consequently no edge of the nonsquarefree graph
can have both endpoints in one state.
The state is constant on a half-open interval of values of $N$, and the
successive intervals form a contiguous partition of $[u,L+1)$.

The reverse data for a candidate $x$ form another contiguous interval
partition, beginning at $x$, the first prefix at which $x$ is present.  On
every interval they point to an anchor state whose occupant is exactly $x$.
Candidates in $U$ point to their own anchor;
candidates in $V$ are indexed by their position in the progression
$18\bmod25$.  For $x\in E$, diagonal completeness is proved separately:
if $x^2+1$ is nonsquarefree, then some odd prime $p\equiv1\pmod4$ satisfies
$p^2\mid x^2+1$, and a tabulated root, verified directly modulo $p^2$, of
$z^2\equiv-1\pmod {p^2}$ places $x$ in the indexed diagonal set.

Fixing $N$ and following these reverse assignments defines a map
$c_N:\mathcal D_N\to\CN$.  The interval-cover conditions show that $c_N$ is
defined for every diagonal candidate.  If $x\neq y$ and $c_N(x)=c_N(y)$,
then they occupy different positions of the same anchor state, so the local
squarefreeness proof gives $xy+1$ squarefree.  Hence $c_N$ is a proper
colouring with colour set $\CN$, exactly as required.
\end{proof}

\section{Uniform tail reductions from five million}

Let
\[
  \DN=\{n\in[1,N]:n\equiv18\pmod {25}\}.
\]
For a Hall completion $H_N(B)$ put
\[
  \RN(B)=H_N(B)\setminus(\CN\cup\DN).
\]
 We first eliminate defects supported only on $\DN$, and then obtain a lower
 bound for the residual set whenever it is nonempty.

\subsection{Pure opposite-base matching}

For each parity, split the source $\DN$ into that parity and the target
$\CN$ into the opposite parity.  The target block has at least the size of
the source block: every $x\in\DN$ satisfies $x\geq18$, and
$x\mapsto x-11$ lies in $[1,N]$, changes the residue $18$ to $7$ modulo $25$,
and reverses parity.  For opposite parities, $xy+1$ is odd, so the
square prime $2$ is inactive; moreover
$xy+1\equiv18\cdot7+1\equiv2\pmod {25}$, so $5$ is inactive.

For a fixed pivot, nonsquarefree values are divided into three disjoint
prime ranges:
\[
 p\leq47,\qquad 47<p\leq N/26,\qquad p>N/26.
\]
 The following three lemmas bound the corresponding exceptional sets.

\begin{lemma}[Small-prime contribution]\label{lem:pure-global-small}
Let $N\geq\FiveMillion$, let $X\subseteq[1,N]$ lie in one residue class
modulo $50$, and let $b$ be any pivot.  Then
\[
 \frac{\#\{x\in X:\exists p\in
 \{3,7,11,13,17,19,23,29,31,37,41,43,47\},\
 p^2\mid bx+1\}}{N}
 \leq \frac{S_0}{\FiveMillion},
\]
where
\[
 S_0=100001\sum_{p}\frac1{p^2}+13
\]
and the sum is over the displayed thirteen primes.
\end{lemma}

\begin{proof}
For a fixed $p$, the conditions $x\bmod50$ and $p^2\mid bx+1$ select at
most one class modulo $50p^2$, hence at most
\[
 \left\lfloor\frac{N+1}{50p^2}\right\rfloor+1
\]
points.  Sum this bound over the thirteen primes.  After division by $N$,
the endpoint term is decreasing for $N\geq\FiveMillion$, so its maximum is
the displayed value at $\FiveMillion$.
\end{proof}

\begin{lemma}[Medium-prime contribution]\label{lem:pure-global-medium}
Under the same hypotheses, the contribution of
$47<p\leq N/26$ is at most
\[
 \frac{3887}{25\cdot10^6}+\frac{59}{10^4}.
\]
\end{lemma}

\begin{proof}
The residue-class union bound gives the reciprocal-square contribution
\[
 \frac1{25}\sum_{47<p\leq N/26}\frac1{p^2}.
\]
The reciprocal-square estimate used here is less than $3887/10^6$.  The
endpoint term contributes $\pi(N/26)/N$, and an explicit prime-count estimate bounds
this by $59/10^4$ for $N\geq\FiveMillion$
\cite{RosserSchoenfeld,DusartPrimes}.  Adding the two terms proves the claim.
\end{proof}

\begin{lemma}[High-prime contribution]\label{lem:pure-global-high}
Assume in addition that $1\leq b\leq N$ and that $bx+1$ is coprime to $10$
for every $x\in X$.  Then the contribution of primes $p>N/26$ is at most
\[
 \frac{11264064/3125}{\FiveMillion}.
\]
\end{lemma}

\begin{proof}
Write $bx+1=p^2m$.  Since $p>N/26$, we have
$p\geq\lfloor N/26\rfloor+1$; while $b,x\leq N$, one has
\[
 bx+1\leq N^2+1
 <677\bigl(\lfloor N/26\rfloor+1\bigr)^2.
\]
Hence $m\leq676$.  After writing the fixed mod-$50$ progression
$x=r+50t$, the equation becomes
\begin{equation}\label{eq:transformed-root}
 p^2m=(br+1)+50bt.
\end{equation}
Let
\[
 S(b)=\{q:q\text{ is prime},\ q\mid b,\ q\notin\{2,5\}\}.
\]
Reducing \eqref{eq:transformed-root} modulo $q\in S(b)$ gives
$p^2m\equiv1\pmod q$.  Thus $q\nmid m$ and $m$ is a nonzero quadratic
residue modulo $q$.  Together with the parity and mod-$5$ condition, these
congruences leave a finite set of possible $m\leq676$.

For fixed $m$, define the root classes to be the residues $s\pmod {50b}$
satisfying
\[
 s^2m\equiv br+1\pmod {50b}.
\]
Every possible prime $p$ lies in one of these classes.  If two solutions in
one class have $p_1<p_2$, then $p_2-p_1\geq50b$.  Subtracting their two
instances of \eqref{eq:transformed-root} gives
\[
 50b(t_2-t_1)=m(p_2-p_1)(p_2+p_1),
\]
and $p_1>N/26$ yields the spacing inequality
\[
 t_2-t_1\geq 2m(N/26).
\]
 Summing the resulting fibre bound over the admissible $m$ and their root
 classes bounds the number of possible $t$.  The finite calculation lists
 exactly those $m$ and root classes satisfying the quadratic-residue
 conditions above, indexed by $S(b)$ rather than by the value of $b$.
 Since the elements of $S(b)$ are distinct prime divisors of $b$,
 $\prod_{q\in S(b)}q\leq b\leq N$; hence every pivot is covered by one of
 these cases.  Applying the spacing inequality in each case gives the
 coefficient $11264064/3125$.  After division by $N$, the resulting bound is
 largest at $N=\FiveMillion$.
\end{proof}

\begin{theorem}[Uniform pure opposite-base matching]
\label{thm:pure-parity-matching}
For every $N\geq\FiveMillion$, the squarefree-edge graph from $\DN$ to
$\CN$ has a matching saturating $\DN$.  Consequently every compatible
$B\subseteq\DN$ satisfies \eqref{eq:hall-defect}.
\end{theorem}

\begin{proof}
 The three contributions satisfy the exact rational inequality
\begin{equation}\label{eq:pure-envelope}
 \frac{S_0}{\FiveMillion}
 +\left(\frac{3887}{25\cdot10^6}+\frac{59}{10^4}\right)
 +\frac{11264064/3125}{\FiveMillion}
 \leq\frac1{100}-\frac1{\FiveMillion}.
\end{equation}
Lemmas~\ref{lem:pure-global-small}--\ref{lem:pure-global-high} therefore
show, in each parity block and after interchanging source and target, that
every vertex has squarefree degree greater than half of the opposite block.
The elementary dense bipartite matching lemma now applies.  Indeed, if a
source subset has at most half the target size, the neighbourhood of any one
of its vertices is already larger than the subset.  If it has more than half
the target size and missed a target vertex, all neighbours of that target
would lie in a source complement of size less than half the source block,
contradicting the reverse degree bound.  Thus Hall's condition holds.  The two
parity matchings combine
because their target parities are disjoint.  Finally, a matching from
$B\subseteq\DN$ into $\CN\setminus\TN(B)$ gives
$|B|\leq|\CN|-|\TN(B)|$.
\end{proof}

\subsection{A quantitative mixed residual}

\begin{proposition}[Uniform mixed degree]\label{prop:uniform-mixed-degree}
Let $N\geq\FiveMillion$, let $B$ be a compatible outside set, and let
$b\in\RN(B)$.  Then $b$ has more than
\[
  \frac N{50}+\frac N{525}
\]
squarefree neighbours in each of $\CN$ and $\DN$.
\end{proposition}

\begin{proof}
Since $\TN(B)\subseteq\CN$, removing $\CN\cup\DN$ from $H_N(B)$ leaves
exactly $B\setminus\DN$.  Thus $b\in B$ and
$b\not\equiv7,18\pmod {25}$.  For $x\equiv7$ or $18\pmod {25}$, the
congruence $25\mid bx+1$ would force, respectively,
$b\equiv7$ or $18\pmod {25}$.  Hence the square $5^2$ cannot divide either
base form associated with a residual pivot.
Split a nonsquarefree form $bx+1$ into the three ranges
\[
 p\in\{2,3,7\},\qquad 7<p\leq N/55,\qquad p>N/55.
\]
The constants
\[
 \tau_7=
 \frac{263529083909042886517376461184337967}
 {8573456796637692379906289787841000000},
 \qquad
 \rho_7=\frac{221926420176}{12755647965025}.
\]
 have the following meanings: $\tau_7$ bounds the reciprocal-square sum for
 primes larger than $7$, and $\rho_7$ bounds the number of admissible roots of
 the transformed equations.
For $p\in\{2,3,7\}$, simultaneous reduction modulo $25p^2$ gives the exact
normalized upper bound $17/1225$.  For $7<p\leq N/55$, the union bound over
the unique relevant residue class gives $\tau_7/25$.

For $p>N/55$, write $x=r+25t$ with $r\in\{7,18\}$ and
$bx+1=p^2m$.  The same comparison with $N^2+1$ now gives $m\leq55^2=3025$.
 There are two transformed equations.  If
$5\nmid br+1$, then $5\nmid m$ and
\[
 p^2m=(br+1)+25bt.
\]
If $5\mid br+1$, then $p\neq5$ forces $5\mid m$; writing $m=5m'$ gives
$m'\leq605$ and, after division by $5$,
\[
 p^2m'=\frac{br+1}{5}+5bt.
\]
 We call these the normal and mod-$5$-twisted equations, respectively.
The nonzero quadratic-residue conditions modulo the prime divisors of $b$,
treated separately in the three even $2$-adic classes and the two odd classes,
give the common upper bound $\rho_7/6$.
The base progression contains at least $N/25-1$ points.  The exact margin is
\[
 \frac{17}{1225}
 +\frac{\tau_7}{25}
 +\frac{\rho_7}6
 +\frac{1+19601/49}{\FiveMillion}
 <\frac1{50}-\frac1{525}.
\]
The endpoint correction decreases with $N$, so subtracting all bad points
from the base progression gives the stated degree in both directions.
\end{proof}

\begin{corollary}[Uniform mixed residual]\label{cor:uniform-mixed-residual}
If $N\geq\FiveMillion$ and a compatible outside set $B$ violates
\eqref{eq:hall-defect}, then
\begin{equation}\label{eq:mixed-residual}
  \frac{2N}{525}-1<|\RN(B)|.
\end{equation}
In particular $|\RN(B)|>19046$, and one of the five valuation parts used
below has more than $3809$ elements.
\end{corollary}

\begin{proof}
Theorem~\ref{thm:pure-parity-matching} implies that $\RN(B)$ is nonempty;
choose $b\in\RN(B)$.  Apply Proposition~\ref{prop:uniform-mixed-degree} to
both base classes.  Let $d_7$ and $d_{18}$ be the resulting squarefree
degrees.  The $\CN$-neighbours are disjoint from $\TN(B)$, and the
$\DN$-neighbours are disjoint from $B\cap\DN$, because products internal to
$B$ are nonsquarefree.  Combining these two facts with the strict defect and
the partition
$B=(B\setminus\DN)\mathbin{\dot\cup}(B\cap\DN)$ gives the exact
degree-subtraction inequality
\[
 d_7+d_{18}<|B\setminus\DN|+|\DN|.
\]
Here $B\setminus\DN=\RN(B)$ and $|\DN|\leq N/25+1$.  Since both degrees are
strictly larger than $N/50+N/525$, we obtain
\[
 \frac{2N}{525}-1<|\RN(B)|.
\]
At $N=\FiveMillion$ the left side is greater than $19046$, and it increases
with $N$.  The five valuation parts are disjoint and cover $\RN(B)$.  If
each had size at most $3809$, their union would have size at most
$5\cdot3809=19045$, a contradiction.
\end{proof}

\section{Valuation classes and pivot estimates}

 The five valuation classes record the $2$-adic information needed by
the congruence counts.  They are
\[
 E_1,\ E_2,\ E_3,\ O_1,\ O_3,
\]
where, respectively,
\[
 x\equiv2\pmod4,\quad x\equiv4\pmod8,\quad x\equiv0\pmod8,
 \quad x\equiv1\pmod4,\quad x\equiv3\pmod4.
\]
Thus $E_1,E_2,E_3$ mean $v_2(x)=1$, $v_2(x)=2$, and $v_2(x)\geq3$,
respectively, while $O_1,O_3$ are the two odd residues modulo $4$.  These five
 conditions are disjoint and exhaustive.
Within a class we next split by the nine residues modulo $9$.  This records
exactly when $3^2$ can divide several pivot forms simultaneously.  When a
finer separation is required, we use fibres modulo $7^2$ or $11^2$ for the
same reason: they control simultaneous divisibility by the next small prime
squares.

\begin{lemma}[Valuation partition]\label{lem:valuation-partition}
The five valuation parts are pairwise disjoint and their union is
$\RN(B)$.  Consequently their cardinalities sum to $|\RN(B)|$.
\end{lemma}

\begin{proof}
 Every even integer satisfies exactly one of $v_2(x)=1$, $v_2(x)=2$, or
 $v_2(x)\geq3$, and every odd integer is congruent to exactly one of $1$ and
 $3$ modulo $4$.  These alternatives are pairwise disjoint and exhaustive,
 so the cardinalities add.
\end{proof}

 The next lemma treats the odd valuation classes in the five-to-ten-million
 range.  It distinguishes a set spread across sufficiently many fibres from a
 concentration that admits a direct cardinality bound.

\begin{lemma}[Nine-cell capacity dichotomy]\label{lem:nine-cell-capacity}
 Fix one odd valuation class after subsets already covered by earlier bounds
 have been removed, and let $\mathcal R$ be the remaining set.  Join a
residue label $c\pmod 9$ to a residue $r\pmod {49}$ when some $x\in\mathcal R$
satisfies $x\equiv c\pmod 9$ and $x\equiv r\pmod {49}$.  For a set $Y$ of
mod-$9$ labels, let $\Gamma_{49}(Y)$ be its neighbourhood in the mod-$49$
residues.  Then either eight distinct mod-$9$ labels can be assigned occupied
mod-$49$ fibres, with no fibre used more than twice, or there is a set $Y$ for
which
\[
  2\,|\Gamma_{49}(Y)|+1<|Y|.
\]
\end{lemma}

\begin{proof}
Apply Hall's theorem to the bipartite graph whose left side is the nine
mod-$9$ labels and whose right side consists of two copies of every mod-$49$
 residue together with one additional vertex.  A saturating matching uses the additional vertex
at most once.  Every other matched edge has, by definition, an occupant in
$\mathcal R$; choosing one produces eight pivots in distinct mod-$9$ cells,
with each mod-$49$ residue used at most twice.  If no saturating matching
exists, Hall supplies $Y$ with the displayed capacity defect.
\end{proof}

 The counting identity used after the pivots are selected begins with the
 disjoint partition
\[
 H_N(B)=\RN(B)\mathbin{\dot\cup}
       \bigl(H_N(B)\cap(\CN\cup\DN)\bigr).
\]
Every element of $H_N(B)$ has nonsquarefree product plus one with every other
element of $H_N(B)$: this follows from compatibility inside $B$, the
definition of $\TN(B)$ for cross-products, and divisibility by $25$ inside
$\TN(B)\subseteq\CN$.  Fix three pivots $P=\{p_1,p_2,p_3\}$ in the completion
and a prime cutoff $z$.  For a base point $x$, let $k(x)$ be the number of
pivots for which $xp_i+1$ has a square-prime divisor at most $z$, and let
$\ell(x)=3-k(x)$ count the remaining large-prime events.  Pointwise,
\[
 1\leq \frac12\mathbf 1_{k(x)=3}
       +\frac12\mathbf 1_{k(x)\geq2}
       +\frac12\ell(x).
\]
Indeed, the three possible ranges $k=3$, $k=2$, and $k\leq1$ make the
right-hand side at least one.  Summing gives the exact three-pivot inequality
\begin{equation}\label{eq:mixed-half-decomposition}
\begin{split}
 |H_N(B)|\leq{}&|\RN(B)|
 +\frac12\#\{x:\text{all three finite-prime events occur}\}\\
 &+\frac12\#\{x:\text{at least two finite-prime events occur}\}\\
 &+\frac12\sum_{p\in P}
       \#\{x:\exists q>z\text{ prime with }q^2\mid px+1\}.
\end{split}
\end{equation}
The points $x$ in the three counted sets are restricted to the base part
$H_N(B)\cap(\CN\cup\DN)$.  The smaller ranges also use four- and six-pivot
versions of this same double-counting principle; from twenty million onward,
all ten cases use exactly \eqref{eq:mixed-half-decomposition}.

For a case $\beta$, let $P_\beta\subseteq\RN(B)$ be its selected pivots.
After division by $N$, the residual bound and the three terms in
\eqref{eq:mixed-half-decomposition} give an inequality of the form
\begin{equation}\label{eq:branch-budget}
 \frac{|H_N(B)|}{N}
 \leq D_\beta+F_\beta+Q_\beta
 <\frac1{25}-\frac7{25L}
 \leq\frac1{25}-\frac7{25N}.
\end{equation}
Here $L$ is the lower endpoint of the interval under consideration.
 The term $D_\beta$ bounds $|\RN(B)|/N$.  It is a diagonal bound because every
$x\in\RN(B)\subseteq H_N(B)$ satisfies that $x^2+1$ is nonsquarefree; the
valuation and mod-$9$ information of the selected case sharpens this count.
 Any valuation classes already known to be small contribute their explicit
 cardinality bounds to $D_\beta$.  The term $F_\beta$ comprises the first two
 finite-prime counts in \eqref{eq:mixed-half-decomposition}, and
 $Q_\beta$ is the large-prime sum, bounded by reciprocal-square estimates and
the transformed-root spacing argument.  The middle inequality
is an exact rational comparison, while the last follows because
$1/25-7/(25N)$ increases with $N$.  Equation
\eqref{eq:branch-budget} therefore contradicts a strict Hall defect.

\begin{theorem}[The five-to-ten-million range]
\label{prop:five-million-interval-close}
The Hall inequality holds for
\[
 \FiveMillion\leq N<\TenMillion.
\]
\end{theorem}

\begin{proof}
Assume a strict defect.  Corollary~\ref{cor:uniform-mixed-residual} gives
 $|\RN(B)|>19046$.  This exceeds the total number removed in any of the case
 reductions below; the largest such total is $128$.  We consider the five
 valuation classes in turn.

First examine $E_1$.  Its mod-$9$ cells and mod-$49$ or mod-$121$ fibres
 either give a three- or four-pivot configuration, or imply $|E_1|\leq10$.
 Under the latter assumption, the nine-cell dichotomy for $E_2$ gives either
 two occupied cells, one sufficiently occupied cell, or $|E_2|\leq10$.
 If both classes satisfy these bounds, $19$ points in $E_3$ give the next
three-pivot configuration; otherwise $|E_3|\leq18$.  Thus all three even
classes together contribute at most
\[
 10+10+18=38
\]
points in the only case that reaches the odd classes.

If one odd class has at most $45$ points, then at most $38+45=83$ points lie
outside the other odd class.  The residual lower bound leaves that other
 class nonempty after its mod-$9$ cells of size at most $5$ have been removed: more than
$128$ residual points and at most $83$ outside the class leave more than $45$
inside it, whereas the nine sparse cells contain at most $9\cdot5=45$ points.
 Hence at least one mod-$9$ cell contains at least six points.  The number of
 such cells is between one and nine.  The corresponding pivot estimate is
 proved separately for each possible number, with a common estimate for six
 through nine cells.
 Otherwise both odd classes contain at least $46$ points.  In each class,
 remove every mod-$9$ cell of size at most $5$; this removes at most
$9\cdot5=45$ points per class.  Together with the even classes, at most
\[
 38+45+45=128
\]
 points have been removed, so a mod-$9$ cell of size at least six remains.
 Either the remaining dense cells contain disjoint triples, giving the
 six-pivot case, or their
union has at most $11$ cells and yields one of the four-pivot cases.  These
 alternatives exhaust the possibilities by the nine-cell capacity
 dichotomy.

 In every alternative the chosen pivots lie in $H_N(B)$ and the
corresponding finite-prime, transformed-root, and diagonal estimates give
\eqref{eq:branch-budget} with $L=\FiveMillion$.  The strict rational
inequality contradicts the assumed defect.
\end{proof}

\begin{theorem}[The ten-to-twenty-million range]
\label{prop:ten-million-interval-close}
The Hall inequality holds for
\[
 \TenMillion\leq N<\TwentyMillion.
\]
\end{theorem}

\begin{proof}
Again $|\RN(B)|>19046$.  Put
\[
 c_N=\left\lfloor\frac{N-1}{1000001}\right\rfloor+1;
\]
for $N<\TwentyMillion$ one has $c_N\leq20$.  More than $c_N$ points in one
mod-$9$ cell contain a pair at distance at most $10^6$.  In $E_1$ and then
$E_2$, the case analysis first asks whether two distinct cells have more than
$c_N$ points, producing two close pairs and hence four pivots.  If not, it
asks whether one cell has more than $c_N$ points, producing a three-pivot
configuration while all other cells have size at most $c_N$.  If neither
happens, the entire valuation class has size at most $9c_N$.  The same
one-cell test is then applied to $E_3$.

 If all three even classes satisfy these bounds, their union has size at most
$27c_N$.  The valuation partition and residual lower bound ensure that at least $91$ residual pivots remain in the two odd classes.  If one odd class
has at most $45$ elements, the other has at least $46$; if neither does, both
 have at least $46$.  The corresponding one-class or two-class argument gives
a three-pivot configuration.  Thus the two-dense-cell alternatives give four
pivots, while the one-cell and odd alternatives give three, and the cases are
exhaustive.

For each of these alternatives, the selected three or four pivots give the
same decomposition as in \eqref{eq:branch-budget}.  Direct residue
 counts give $F_\beta$, the prime-square fibres give $Q_\beta$, and the bounds
 for the remaining cells are included in $D_\beta$.  Substitution puts the total
strictly below $1/25-7/(25N)$, so a strict defect is impossible.
\end{proof}

\subsection{A ten-case decomposition from twenty million onward}

The later ranges use the following exhaustive case decomposition.  Put
\[
 g_N=\left\lfloor\frac{N-1}{20001}\right\rfloor+1,
 \qquad
 \delta=\frac1{20001}+\frac1{\TwentyMillion}.
\]
The number $g_N$ is an exact cardinality cap, and
$g_N/N\leq\delta$ for $N\geq\TwentyMillion$.  A class containing more than
$g_N$ points has two points at distance at most $20000$ and, since
$g_N\geq2$, a third distinct point.  Call the resulting triple
\emph{generic} when its three members are not all congruent modulo $9$.
Otherwise, either the whole valuation class lies in that common mod-$9$ cell,
or a point outside the cell can replace the third member and make the triple
 generic.  Hence, if no generic triple exists, the entire valuation class lies
 in one residue class modulo $9$.

\begin{lemma}[Ten-case decomposition]\label{lem:global-ten-branch-allocation}
Let $N\geq\TwentyMillion$ and suppose that $B$ is a compatible outside set
with a strict Hall defect.  Then one of the ten cases in
Table~\ref{tab:ten-branch-structure} applies.
\end{lemma}

\begin{proof}
The uniform cutoff-$19$ degree theorem strengthens the mixed-residual estimate
to $0.0045N<|\RN(B)|$.  Inspect $E_1,E_2,E_3$ in order.  If the current class
has more than $g_N$ members, the preceding close-triple and mod-$9$
 dichotomy gives one of its two cases.  Otherwise the class has cardinality at
 most $g_N$.  If all three even classes satisfy this bound, the remaining
 residual elements lie in the two odd classes.  When both odd classes are
 nonempty, the pigeonhole principle on intervals of length $460$ gives a pair
 at distance at most $459$; when only one odd class remains, intervals of
 length $230$ give a pair at distance at most $229$.  Adding a third point and
 applying the same mod-$9$ dichotomy gives the final four cases.  These
 alternatives exhaust all five valuation classes.
\end{proof}

The eight symbols in the last column of Table~\ref{tab:ten-branch-structure}
are the eight diagonal subsets that are counted separately: all residual
points; one common $E_1$ cell; the low-$2$-adic classes; one common $E_2$
cell; the union of the two odd classes; that union with one common cell; one
 odd class; and one odd class with a common cell.  The entry $j\delta$ accounts
 for at most $j$ even classes satisfying the preceding size bounds.  In the
 formalization these conditions and bounds are represented by
 \texttt{TwentyMillionBranchApplies} and
\texttt{twentyMillionBranchResidualPayment}.

\begin{table}[ht]
\centering
\footnotesize
\renewcommand{\arraystretch}{1.12}
\begin{tabularx}{\textwidth}{@{}>{\raggedright\arraybackslash}p{0.19\textwidth}Y Y@{}}
\toprule
 Case & Condition after the preceding size bounds hold & Set used to bound $D_\beta$\\
\midrule
$E_1$ generic & close triple in $E_1$, not constant mod $9$ & all residual points\\
$E_1$ common & all of $E_1$ lies in one mod-$9$ cell & common $E_1$ cell\\
$E_2$ generic & $|E_1|\leq g_N$; generic triple in $E_2$ & all residual points\\
$E_2$ common & $|E_1|\leq g_N$; all of $E_2$ lies in one cell & common $E_2$ cell $+\delta$\\
$E_3$ generic & $|E_1|,|E_2|\leq g_N$; generic triple in $E_3$ & $E_1\cup E_2$, plus $2\delta$\\
$E_3$ common & same, with all of $E_3$ in one cell & $E_1\cup E_2$, plus $2\delta$\\
two odd, generic & all even classes have size $\leq g_N$; generic odd triple & $O_1\cup O_3$, plus $3\delta$\\
two odd, common & same, with the relevant odd class in one cell & $O_1\cup O_3$ and that cell, plus $3\delta$\\
one odd, generic & all even classes are small; the other odd class is empty & remaining odd class, plus $3\delta$\\
one odd, common & same, with the remaining odd class in one cell & that common cell, plus $3\delta$\\
\bottomrule
\end{tabularx}
\caption{The ten exhaustive cases used for $N\geq\TwentyMillion$.}
\label{tab:ten-branch-structure}
\end{table}

\begin{theorem}[The twenty-to-forty-million range]
\label{prop:twenty-million-interval-close}
The Hall inequality holds for
\[
 \TwentyMillion\leq N<\FortyMillion.
\]
\end{theorem}

\begin{proof}
 Lemma~\ref{lem:global-ten-branch-allocation} gives one of the ten cases in
 Table~\ref{tab:ten-branch-structure}.
 The six even cases use finite cutoff $23$; the four odd cases use cutoff $19$.
 For the pivots selected in case $\beta$, the decomposition above gives
\[
 \frac{|H_N(B)|}{N}\leq T_\beta,
\]
where $T_\beta=D_\beta+F_\beta+Q_\beta$.  We display one complete comparison
so that the normalization and strict margin can be checked without consulting
the formal source.  Consider the case in which $E_1$ has at most
 $g_N=\lfloor(N-1)/20001\rfloor+1$ elements and $E_2$ contains a triple that
is not contained in one residue class modulo $9$.  Put
\[
 \tau_{23}=
 \frac{64081802747648035629863}
 {7596668444022826249000000}.
\]
 This is the reciprocal-square upper bound for primes larger than $23$.  In
 this case the diagonal estimate for all residual points already includes
 $E_1$, so no separate $\delta$ term is needed.  The three upper bounds
in \eqref{eq:branch-budget} are
\[
 D_\beta=\frac{25289550}{10^9},\qquad
 F_\beta=\frac{8685}{10^6},\qquad
 Q_\beta=\frac{3\tau_{23}}{25}
          +\frac12\frac{10006474}{10^9}.
\]
Consequently
\begin{align*}
 T_\beta
 &=\frac{303791140850460908947114523}
 {7596668444022826249000000000},\\
 \left(\frac1{25}-\frac7{25\cdot\TwentyMillion}\right)-T_\beta
 &=\frac{75490557093924693317991}
 {7596668444022826249000000000}>0.
\end{align*}
 The other nine cases differ only in the valuation classes already bounded, the
diagonal subset from Table~\ref{tab:ten-branch-structure}, and the finite- and
large-prime bounds.  Substituting their exact rational values proves, in every
case,
\[
 T_\beta<\frac1{25}-\frac7{25\cdot\TwentyMillion}
 \leq\frac1{25}-\frac7{25N}
\]
and hence contradicts the assumed Hall defect.
\end{proof}

\section{From forty million to two billion}

\begin{theorem}[The forty-to-two-hundred-million range]
\label{thm:mixed-joint-tail-close}
The Hall inequality holds for
\[
 \FortyMillion\leq N<\TwoHundredMillion.
\]
\end{theorem}

\begin{proof}
The interval is divided into the six half-open blocks with endpoints
\[
 40,\ 50,\ 70,\ 80,\ 100,\ 150,\ 200
\quad\text{million}.
\]
 Lemma~\ref{lem:global-ten-branch-allocation} gives one of the ten cases.  On each
block, $D_\beta$ is the corresponding diagonal upper bound plus the exact
cardinality bounds for preceding small classes.  The finite-prime terms
$F_\beta$, in the order of Table~\ref{tab:ten-branch-structure}, are
\[
 10^{-6}(8685,12616,8685,12616,8685,12616,
 19420,20878,26643,29459).
\]
The term $Q_\beta$ uses cutoff $23$ in the six even cases and cutoff $19$
in the four odd cases, together with the transformed-root bound for the
current block.  These are uniform inequalities for every $N$ in the block:
 the residue counts and fibre-spacing theorem give the first inequality in
 \eqref{eq:branch-budget} uniformly throughout the block.  The
final rational comparison is made at the lower endpoint $L$, after which the
monotonicity of $1/25-7/(25N)$ gives the result throughout the block.  Hence
no strict defect exists.
\end{proof}

\begin{theorem}[The two-hundred-million-to-two-billion range]
\label{thm:actual-support-mixed-close}
The Hall inequality holds for
\[
 \TwoHundredMillion\leq N<\TwoBillion.
\]
\end{theorem}

\begin{proof}
The four subintervals have endpoints
\[
 200,\ 300,\ 500,\ 1000,\ 2000
\quad\text{million}
\]
and transformed-root splits $75,85,100,125$, respectively.  Thus the
large-prime threshold is $p>N/s$ with $s$ equal to the listed split, and the
complementary quotient is at most $s^2$.  The normal and mod-$5$-twisted
cases use, respectively, the integers in $[1,s^2]$---or the quotients after
dividing by $5$---that satisfy the nonzero quadratic-residue conditions
forced by the prime divisors of the pivot.  The fibre-spacing argument following
 \eqref{eq:transformed-root} converts their exact cardinalities into a root
 bound for each subinterval.

 Lemma~\ref{lem:global-ten-branch-allocation} gives one of the same ten
 cases.  For a subinterval $\rho$ and case $\beta$, define
\[
 T_{\rho,\beta}
 =D_{\rho,\beta}+F_\beta+Q_{\rho,\beta},
\]
where $D_{\rho,\beta}$ is the bound for the relevant one of the eight
diagonal subsets described before Table~\ref{tab:ten-branch-structure},
$F_\beta$ is the finite-prime vector displayed in the preceding proof, and
$Q_{\rho,\beta}$ uses the subinterval's transformed-root upper bound
\[
 10^{-9}(9000000,\ 7700000,\ 6400000,\ 5000000).
\]
 The finite sets consist exactly of the complementary quotients satisfying
 the displayed congruence conditions.  The equation $p^2m=bx+1$ and the prime
 divisors of $b$ show that every possible quotient belongs to the appropriate
 set; the spacing lemma then gives its contribution to the cardinality bound.
 Combining this with the
diagonal and finite-prime estimates proves
$|H_N(B)|/N\leq T_{\rho,\beta}$, and exact rational
arithmetic gives
\[
 T_{\rho,\beta}
 <\frac1{25}-\frac7{25L_\rho}
 \leq\frac1{25}-\frac7{25N}
\]
 for every subinterval and case.  This contradicts a strict defect.
\end{proof}

\section{The high tail}

For the high tail the same ten mathematical cases remain valid.  Only the
 uniform bounds for the relevant diagonal subset and for the admissible roots
 change with $N$, so the finite part is covered by consecutive
intervals and the remaining part by one monotone argument.

\begin{theorem}[Finite high tail]\label{thm:finite-high-tail}
The Hall inequality holds for
\[
 \TwoBillion\leq N<\FiveHundredBillion.
\]
\end{theorem}

\begin{proof}
The half-open row endpoints are
\[
 2,\ 2.3,\ 2.8,\ 4.2,\ 6,\ 12,\ 500
\quad\text{billion}.
\]
The endpoints are contiguous, so every $N$ in the displayed range belongs
to one half-open interval.  On that interval, the finite-prime cutoff fixes
$F_\beta$, while the nonzero quadratic-residue restrictions and
 \eqref{eq:transformed-root} bound $Q_\beta$.  Every solution is first shown to
 belong to the corresponding finite list, after which the spacing theorem
 bounds each root fibre.  Lemma
~\ref{lem:global-ten-branch-allocation} selects a case, and exact rational
comparison puts $D_\beta+F_\beta+Q_\beta$ strictly below
$1/25-7/(25N)$.  Thus \eqref{eq:hall-defect} follows.
\end{proof}

\begin{theorem}[Unbounded high tail]\label{thm:unbounded-high-tail}
The Hall inequality holds for every
\[
 N\geq\FiveHundredBillion.
\]
\end{theorem}

\begin{proof}
Use split $120$, so large prime squares are those with $p>N/120$ and their
 complementary quotients satisfy $m\leq120^2$.  At the left endpoint,
\[
 840^4\leq\FiveHundredBillion\leq N.
\]
Consequently the factor
$\bigl(\sqrt{\sqrt N}+1\bigr)/N$ in the root count is at most
 $841/840^4$.  In the formalization, the integer $840$ is stored in the field
 \texttt{rootFloor}; its sole role is to replace the variable fourth-root
 factor by this fixed rational bound.
 An explicit sieve over primes up to $47$, together with the quadratic-residue
 conditions and fibre spacing, gives the common transformed-root bound
\[
 \frac{8062340}{10^9}.
\]
 The eight diagonal upper bounds are
\[
 10^{-9}(26005710,20944421,17208292,16649746,
 13003280,7437057,6604680,917036),
\]
The diagonal count for each of the eight subsets is periodic after imposing
 the selected small-prime congruences.  Grouping those congruences by their
 counting moduli gives
\[
 (1821,2892,2627,2295,3643,4590,3643,3643).
\]
For each diagonal subset, the relevant residue classes are counted modulo its
 displayed modulus; the associated reciprocal-square tail and endpoint
term give the corresponding entry in the diagonal vector.  The split at
$120$, together with the sieve over primes up to $47$ and the spacing inequality after
\eqref{eq:transformed-root}, gives the common root bound.  Substitution in the
ten case formulas puts every total below the target at
$\FiveHundredBillion$.  The prime-count, reciprocal-square, fibre-spacing,
and endpoint terms are all bounded by functions that do not increase after
this endpoint.  Therefore the same inequalities apply for every larger $N$
and rule out a strict Hall defect.
\end{proof}

\begin{theorem}[The range from two billion onward]\label{thm:formal-high-tail-close}
The Hall inequality, and hence the exact upper bound, holds for every
$N\geq\TwoBillion$.
\end{theorem}

\begin{proof}
Theorem~\ref{thm:finite-high-tail} covers
$\TwoBillion\leq N<\FiveHundredBillion$, and
Theorem~\ref{thm:unbounded-high-tail} covers
$N\geq\FiveHundredBillion$.
\end{proof}

\section{Completion of the proof}

\begin{theorem}[The range from forty million onward]
\label{thm:forty-million-tail-assembly}
The Hall inequality, and hence the exact upper bound, holds for every
$N\geq\FortyMillion$.
\end{theorem}

\begin{proof}
Theorems~\ref{thm:mixed-joint-tail-close},
\ref{thm:actual-support-mixed-close}, and
\ref{thm:formal-high-tail-close} cover, respectively,
\[
 [\FortyMillion,\TwoHundredMillion),\qquad
 [\TwoHundredMillion,\TwoBillion),\qquad
 [\TwoBillion,\infty).
\]
\end{proof}

\begin{theorem}[The range from five million onward]
\label{thm:mixed-cell-five-million-close}
The Hall inequality, and hence the exact upper bound, holds for every
$N\geq\FiveMillion$.
\end{theorem}

\begin{proof}
Theorems~\ref{prop:five-million-interval-close},
\ref{prop:ten-million-interval-close}, and
\ref{prop:twenty-million-interval-close} cover
\[
 [\FiveMillion,\TenMillion),\quad
 [\TenMillion,\TwentyMillion),\quad
 [\TwentyMillion,\FortyMillion),
\]
and Theorem~\ref{thm:forty-million-tail-assembly} covers the remaining
unbounded interval.  These four half-open intervals are pairwise disjoint and
cover $[\FiveMillion,\infty)$.
\end{proof}

\begin{theorem}[Erd\H{o}s Problem 848]
\label{thm:erdos-848-complete}
For every integer $N\geq1$, the maximum cardinality of a set
$A\subseteq[1,N]$ for which $ab+1$ is nonsquarefree for every $a,b\in A$ is
\[
 \left|\{1\leq n\leq N:n\equiv7\pmod {25}\}\right|.
\]
Equivalently, the Hall inequality \eqref{eq:hall-defect} holds for every $N$.
\end{theorem}

\begin{proof}
Theorem~\ref{thm:finite-prefix-colouring-close} gives the upper bound for
$N\leq\FiveMillion$, and
Theorem~\ref{thm:mixed-cell-five-million-close} gives it for
 $N\geq\FiveMillion$.  Proposition~\ref{prop:sharpness} proves sharpness.
The Hall formulation is equivalent by
Theorem~\ref{thm:hall-equivalence}.
\end{proof}

\section{Correspondence with the Lean formalization}
\label{sec:lean-correspondence}

 This section provides a dictionary between the mathematical proof above and
 the formal source.  Each name below denotes a definition or theorem in the
 dependency chain of the final declaration.

\subsubsection*{The problem and the Hall sets}
In \path{lean4/Erdos848/ProblemCore.lean},
\path{NonSquarefreeProductProp} is the quantified condition on all ordered
pairs, including the diagonal, and \path{OriginalA7} is $\CN$.  In
\path{lean4/Erdos848/HallReduction.lean}, \path{OriginalA18} is $\DN$,
\path{hallNonNeighbours} is $\TN(B)$, and \path{hallCompletion} is $H_N(B)$.
The formal hypothesis corresponding to a compatible outside set is the
conjunction of \path{Erdos848OutsideSet N B}, which says that $B$ lies in
$[1,N]\setminus\CN$, and \path{NonSquarefreeProductProp B}, which gives
compatibility inside $B$.  Finally,
\path{lean4/Erdos848/HallPartition.lean} defines \path{hallResidual} and
\path{hallBasePart}; these are exactly $\RN(B)$ and
$H_N(B)\cap(\CN\cup\DN)$.  The theorem
\path{hallCompletion_card_partition} establishes the partition used in
\eqref{eq:mixed-half-decomposition}.

\subsubsection*{The finite prefix}
The structure \path{PrefixColouringState} in
\path{lean4/Erdos848/ProblemCore.lean} has the two fields used in the proof of
Theorem~\ref{thm:finite-prefix-colouring-close}: every diagonal candidate is
sent to \path{OriginalA7 N}, and two distinct candidates with the same colour
have squarefree product plus one.  The interval histories produce this
structure through \path{prefixColouringCertificate} in
\path{lean4/Erdos848/LowRangePrefixTraceChecker.lean}.
 The declaration
\begin{center}\small
\path{GeneratedFiveMillionPrefixTrace.closeThroughFiveMillion}
\end{center}
is used by \path{erdos848_through_five_million}, together with the general
implication from a colouring to the extremal bound,
\path{originalProblem_of_prefixColouringState}.

\subsubsection*{Root counting and the mixed residual}
The congruence imposed on the quotient $m$ is proved by
\path{quotient_modEq_square_of_prime_dvd_pivot} in
\path{lean4/Erdos848/TailFiveMillionActualSupportCore.lean}.  The spacing
inequality following \eqref{eq:transformed-root} is
 \path{transformedEquation_parameter_gap} in
 \path{lean4/Erdos848/TailTransformedRootCounting.lean}.  These lemmas are used
 in the pure high-prime count and in the later normal and mod-$5$-twisted root
 counts.  The two degree bounds of
Proposition~\ref{prop:uniform-mixed-degree} are
 \path{globalMixedA7_degree_lower} and \path{globalMixedA18_degree_lower}; they
 are combined in \path{globalMixedHallResidual_cast_lower_of_defect} to give
\eqref{eq:mixed-residual}.

\subsubsection*{The pivot inequality and the ten cases}
The pointwise counting argument in
\eqref{eq:mixed-half-decomposition} is formalized first as the abstract lemma
\path{card_le_half_three_two_add_half_tail_sum}, and then for the actual Hall
sets as
\begin{center}\small
\path{hallCompletion_card_le_fiveMillionR263MixedHalfComponents}.
\end{center}
Both occur in
\path{lean4/Erdos848/TailR263MixedHalfDecomposition.lean}.  In the latter
 theorem the three summands after \path{hallResidual} are precisely the
 finite-prime and large-prime contributions denoted by $F_\beta$ and
$Q_\beta$ above.  From twenty million onward,
\path{TwentyMillionBranchApplies} states the ten mutually exhaustive
alternatives in Table~\ref{tab:ten-branch-structure}, while
\path{twentyMillionBranchResidualPayment} assigns the corresponding
$D_\beta$.  The theorem
\path{TwentyMillionKernelTerminalCertificate.completion_ratio_le_branchTotal}
 instantiates the displayed Hall-completion inequality and combines its three
 terms with the bounds for the selected case.  This is the formal counterpart of
\eqref{eq:branch-budget}.

\subsubsection*{Interval theorems}
For each interval, the following implication appears in the Lean sources:
\begin{itemize}\small
\item $[\FiveMillion,\TenMillion)$:\par
\path{fiveMillionR263BranchExhaustion_kernel}\par
$\longrightarrow$ \path{erdos848FiveToTenMillionClose}.
\item $[\TenMillion,\TwentyMillion)$:\par
\path{tenMillionKernelBranchExhaustion}\par
$\longrightarrow$ \path{erdos848TenToTwentyMillionClose}.
\item $[\TwentyMillion,\FortyMillion)$:\par
\path{exists_twentyMillionBranchApplies_global19}\par
$\longrightarrow$ \path{erdos848TwentyMillionClose_kernel}.
\item $[\FortyMillion,\TwoHundredMillion)$:\par
\path{fortyMillionTenBranchTotal_lt_target}\par
$\longrightarrow$ \path{erdos848FortyMillionClose_kernel}.
\item $[\TwoHundredMillion,\TwoBillion)$:\par
\path{hybridTenBranchTotal_lt_target}\par
$\longrightarrow$
\path{erdos848PaperTwoHundredMillionToTwoBillionClose_of_certificates}.
\item $[\TwoBillion,\FiveHundredBillion)$:\par
\path{highQrFinalFiniteCertificate}\par
$\longrightarrow$ \path{HighQrFiniteCloseCertificate.close}.
\item $[\FiveHundredBillion,\infty)$:\par
\path{highUnboundedRootTerminal}\par
$\longrightarrow$ \path{erdos848HighUnboundedClose}.
\end{itemize}
These declarations cover consecutive intervals and are combined by the
declaration
\begin{center}\small
\path{Erdos848.PaperGeneratedCertificateProvider.all_N}.
\end{center}

\section{Kernel verification and reproduction}
\label{sec:machine-proof}

The proof has been formalized in Lean~4 with Mathlib
\cite{Lean4,Mathlib}.  The final upper-bound declaration is
\begin{center}\small
\texttt{Erdos848.PaperGeneratedCertificateProvider.all\_N}.
\end{center}
Together with \texttt{originalA7\_has\_property} and
\texttt{erdos848HallStatement\_iff\_originalProblem}, it yields the sharpness
and Hall-equivalence clauses of Theorem~\ref{thm:erdos-848-complete}.
The formal statements use the same diagonal convention, mod-$25$ base
 classes, Hall completion, valuation classes, case conditions, and rational
 bounds as the manuscript.  The repository includes a correspondence table
 recording the Lean declarations associated with each displayed definition
 and inequality.  In particular, the calculation in the proof of
 Theorem~\ref{prop:twenty-million-interval-close} is linked to the
 \texttt{evenTwoGeneric} case.

Programs used to find finite witnesses or candidate quotient sets lie outside the
logical argument.  Their outputs are read as ordinary Lean data, and general
theorems prove the interval coverage, congruence conditions, fibre bounds,
and rational inequalities needed above.  Only the resulting proof terms are
accepted by the Lean kernel.  A \texttt{--trust=0} axiom audit of the final
theorem reports
\[
 \texttt{[propext, Classical.choice, Quot.sound]},
\]
the standard classical and quotient dependencies used by Lean and Mathlib.

The version-independent Zenodo concept DOI for the code archive is
\href{https://doi.org/10.5281/zenodo.21750213}%
{\nolinkurl{10.5281/zenodo.21750213}}; it resolves to the latest archived
 code release.  The repository pins the Lean toolchain and Mathlib dependency
 graph.  The archived proof can be checked with
\begin{quote}\scriptsize\ttfamily
python -B scripts/check\_proof\_state.py
  --require-release-ready --audit-sources
\end{quote}
 and the correspondence, arithmetic, and reference checks can be run with
\begin{quote}\scriptsize\ttfamily
python -B scripts/verify\_paper\_math\_implementation.py\\
python -B scripts/verify\_paper\_lean\_correspondence.py\\
python -B scripts/verify\_paper\_lean\_numbers.py\\
python -B scripts/verify\_reference\_evidence.py\\
\hspace*{2em}--require-cited-coverage --require-entry-checks
\end{quote}

\section*{Disclosure and verification statement}

Multiple artificial-intelligence tools and an author-developed research system
were used in the research and preparation of this article.  The author is
responsible for the mathematical content and conclusions of the article.
All results can also be independently verified through the accompanying
kernel-checked machine-proof project; their formal validity is therefore
subject to reproducible kernel verification rather than the author's assertion
alone.  Further details concerning the systems and their use will be disclosed
separately.

\section*{Acknowledgements}

The author thanks the developers of Lean and Mathlib for the formalization
environment.

\bibliographystyle{plain}
\bibliography{references}

\end{document}
